Verwenden Sie die Methode von Rothstein-Trager, um die folgenden Integrale
auszurechnen.
\begin{teilaufgaben}
\item $\displaystyle
\int \frac{x+1}{x^2+1}\,dx$
\item $\displaystyle
\int \frac{x}{x^2+x+1}\,dx$
\end{teilaufgaben}

\begin{loesung}
\begin{teilaufgaben}
\item Der Nenner $b(x)=x^2+1$ ist bereits quadratfrei und der Zähler $a(x)=x+1$
hat kleineren Grad als der Nenner, die Methode von Rothstein-Trager ist
ohne weitere Vorarbeit anwendbar.
Es sind nichttriviale gemeinsame Teiler von
\begin{align*}
a(x)-cb'(x)
&=x+1-c(2x)
=1+(1-2c)x
=
(1-2c)\biggl(x+\frac{1}{1-2c}\biggr)
\intertext{und}
b(x)&=x^2+1
\end{align*}
zu finden.
Da $b(x)=(x+i)(x-i)$ sind solche Teiler zu finden für
\[
\frac{1}{1-2c}=\pm i
\quad\Rightarrow\quad
1-2c=\mp i
\quad\Rightarrow\quad
2c=1\pm i
\quad\Rightarrow\quad
c=\frac12\pm\frac{i}{2}
\]
und der zugehörige grösste gemeinsame Teiler ist $(x\pm i)$.
Somit ist die Stammfunktion
\begin{align*}
\int \frac{x+1}{x^2+1}\,dx
&=
\biggl(\frac12+\frac{i}2\biggr)\log(x+i)
+
\biggl(\frac12-\frac{i}2\biggr)\log(x-i)
\\
&=
\frac12 \log((x+i)(x-i))
+
\frac{i}2 \log\frac{x+i}{x-i}
\\
&=
\frac12\log(x^2+1)
+\arctan x.
\end{align*}
\item
Mit $a(x)=x$ und $b(x)=x^2+x+1$ sind gemeinsame Teiler von
\[
a(x)-cb'(x)
=
x-2cx-c
=
(1-2c)x-c
\quad\text{und}\quad
b(x)= x^2+x+1
\]
zu finden.
Dazu führen wir den euklidischen Algorithmus für die Polynome $b_0(x)=x^2+x+1$
und $b_1(x)=x-d$ durch, wobei $d=c/(1-2c)$.
Division von $b_0(x)$ durch $b_1(x)$ ergibt
\[
b_0(x) = (x+d+1) b_1(x) + d^2 +d + 1
\]
Der grösste gemeinsame Teiler ist also $b_1(x)$, wenn der Rest
$d^2+d+1$ verschwindet.
Setzt man $d=c/(1-2c)$ für $d$ ein, erhält man die Bedingung
\[
\frac{c^2}{(1-2c)^2}+\frac{c}{1-2c}+1=0
\quad\Rightarrow\quad
\frac{3c^2-3c+1}{4c^2-4c+1}=0.
\]
Die möglichen Werte von $c$ sind daher
\begin{align*}
c_{1,2}
&=
\frac{3\pm\!\sqrt{9-4\cdot 3}}{2\cdot 3}
=
\frac{3\pm\!\sqrt{-3}}{6}
=
\frac{3\pm i\!\sqrt{3}}{6}
\intertext{und die zugehörigen Werte von $d$}
d_{1,2}
&=
\frac{c_{1,2}}{1-2c_{1,2}}
=
\frac{3\pm i \!\sqrt{3}}{6}
\biggl(
1-\frac{3\pm i\!\sqrt{3}}{3}
\biggr)^{-1}
=
\frac{3\pm i \!\sqrt{3}}{6}
\frac{3}{\mp i\!\sqrt{3}}
=
\pm
i
\frac{\!\sqrt{3}}{2}
-\frac12.
\end{align*}
Damit kann man jetzt das Integral zusammensetzen:
\begin{align*}
\int \frac{x}{x^2+x+1}\,dx
&=
c_1\log v_1(x)
+
c_2\log v_2(x)
\\
&=
c_1\log (x-d_1)
+
c_2\log (x-d_2)
\\
&=
\frac{3+i\!\sqrt{3}}{6}\log \biggl(x+\frac12-i\frac{\!\sqrt{3}}2\biggr)
+
\frac{3-i\!\sqrt{3}}{6}\log \biggl(x+\frac12+i\frac{\!\sqrt{3}}2\biggr).
\intertext{Wir trennen Real- und Imaginärteil der Koeffizienten und erhalten}
&=
\frac12
\log\biggl(
\biggl(x+\frac12-i\frac{\!\sqrt{3}}{2}\biggr)
\biggl(x+\frac12+i\frac{\!\sqrt{3}}{2}\biggr)
\biggr)
+\frac{1}{\!\sqrt{3}}
\frac{i}{2}
\log
\frac{
x+\frac12+i\frac{\!\sqrt{3}}{2}
}{
x+\frac12-i\frac{\!\sqrt{3}}{2}
}
\\
&=
\frac12\log \biggl(\biggl(x+\frac12\biggr)^2+\frac34\biggr)
+
\frac{1}{\!\sqrt{3}}
\frac{i}{2}\log
\frac{
2x+1+i\!\sqrt{3}
}{
2x+1-i\!\sqrt{3}
}
\\
&=
\frac12\log \biggl(x^2+x+\frac14+\frac34\biggr)
+
\frac{1}{\!\sqrt{3}}
\frac{i}{2}\log
\frac{
(2x+1)/\!\sqrt{3}+i
}{
(2x+1)/\!\sqrt{3}-i
}
\\
&=
\frac12\log(x^2+x+1) -\frac{1}{\!\sqrt{3}}\arctan\frac{2x+1}{\!\sqrt{3}}.
\qedhere
\end{align*}
\end{teilaufgaben}
\end{loesung}

